\section{Background}

Performance in modern computing systems depends on several interacting factors, including algorithmic complexity, processor architecture, and memory hierarchy.

Algorithmic complexity describes how the runtime of an algorithm grows as input size increases. Algorithms with quadratic complexity scale dramatically worse than algorithms with near-linear growth rates, particularly for large inputs.

Modern processors rely heavily on hierarchical memory systems consisting of multiple levels of cache and main memory \cite{lam1991}. Accessing data stored in cache is significantly faster than accessing data from main memory. As a result, programs that access memory sequentially often perform substantially better than programs that access memory in irregular patterns \cite{mckee2004}.

Additionally, performance scaling in modern computing is constrained by fundamental architectural limits such as Amdahl's law and the so-called memory wall \cite{hill2008, wulf1995}. These constraints suggest that efficient algorithms and memory-aware implementations remain critical even as raw computational power increases.
